\documentclass{article}
\usepackage[T1]{fontenc}
\usepackage{lmodern}
\usepackage{graphicx}
\usepackage{amsmath,amssymb}
\usepackage[a4paper,margin=2cm]{geometry}
\usepackage[absolute,overlay]{textpos}
\setlength{\TPHorizModule}{1mm}
\setlength{\TPVertModule}{1mm}
\usepackage[indonesian]{babel}
\usepackage{hyperref}
\setlength{\emergencystretch}{2em}
% unit: Halaman 46

\begin{document}

% === Halaman 46 (python/native) ===

Perbandingan Elgamal dengan Elliptic Curve Elgamal

• Cipherteks pada EC-Elgamal adalah pasangan  titik 

• PC = ( (kB), (PM

 + kPB) )            (ket: Pb = kunci publik Bob)

• Cipherteks pada Elgamal juga pasangan nilai:

• C = (a, b) = (gk mod p, myB

k mod p)    (ket: yb = kunci publik Bob)

-----------------------------------------------------------------------------------------------------------

• Pada EC\_Elgamal, Bob mengurangkan titik kedua  dari PC dengan hasil kali b  (kB)

• (PM + kPB) – [b(kB)] = PM + k(bB) – b(kB) = PM

• Pada El Gamal, Bob membagi nilai kedua dengan nilai pertama yang dipangkatkan 

dengan kunci privat Bob

• myB

k / (gk)b = mgk*b / gk*b = m

*) Sumber bahan: Debdeep Mukhopadhyay, Elliptic Curve Cryptography ,

 Dept of Computer Sc and Engg IIT Madras

46

\end{document}
